\subsection{Exact zero-knowledge parity proofs from bit commitments}\label{sec:realize:parity}

\begin{definition}[Parity-compatible integer openings]\label{def:realize:candidates}
For $m\geq0$ and a public desired parity
$b\in\{0,1\}$, define
\[
 \mathcal S_{m,b}:=\{t\in\{0,\ldots,m\}:t\equiv b
                                  \pmod 2\}.
\]
If $\mathcal S_{m,b}$ is empty, the verifier
rejects the parity statement deterministically.
Assume $P>m$, so all candidates map injectively
into $\mathbb Z_P$.
\end{definition}

\begin{lemma}[Correct integer sum lies in the permitted set]\label{lem:realize:correct-candidate}
Let $b_1,\ldots,b_m\in\{0,1\}$.
If $\sum_i b_i=b$ in $\F_2$,
then the ordinary integer
$S=\sum_i b_i$ belongs to $\mathcal S_{m,b}$.
\end{lemma}
\begin{proof}
As an integer, $0\leq S\leq m$.
Reduction of the sum modulo $2$
is exactly addition of the bits in $\F_2$,
so $S\equiv b\pmod2$, giving membership.
\end{proof}

\begin{definition}[Concrete parity proof]\label{def:realize:parity-proof}
Given public bit commitments
$C_i=g^{b_i}h^{r_i}$, form $C_T=\prod_iC_i$.
The prover proves:
\begin{enumerate}[label=\textup{(\roman*)}]
\item For each $C_i$ it knows a bit opening,
using Definition~\ref{def:realize:or}
with candidate set $\{0,1\}$.
\item It knows an opening of $C_T$ to an
integer $t\in\mathcal S_{m,b}$, using
Definition~\ref{def:realize:or} with
candidate set $\mathcal S_{m,b}$.
\end{enumerate}
All OR-proofs have independent fresh prover
randomness and independent uniform sigma
verifier challenges. The first bitness
proofs can be executed when the commitments
are bound; the aggregate parity proof is
performed when the public row coefficients
and target parity become known.
\end{definition}

\begin{theorem}[Perfect completeness of the parity compiler]\label{thm:realize:parity-complete}
For honest bit commitments and bits satisfying
$\sum_i b_i=b$ in $\F_2$, the concrete
parity proof accepts with probability $1$.
\end{theorem}
\begin{proof}
Each $C_i$ has the known bit opening
$(b_i,r_i)$, so its two-branch OR proof
is perfectly complete by
Lemma~\ref{lem:realize:or-complete}.
Lemma~\ref{lem:realize:correct-candidate}
places the true integer sum
$S=\sum_i b_i$ in $\mathcal S_{m,b}$.
Lemma~\ref{lem:realize:aggregate}
gives the known aggregate opening
$C_T=g^S h^{R_T}$ with
$R_T=\sum_i r_i$.
The aggregate OR proof is therefore
complete by the same lemma.
All checks pass.
\end{proof}

\begin{lemma}[An incorrect parity gives two different aggregate openings]\label{lem:realize:parity-equivocation}
Suppose each $C_i$ has an extracted
valid bit opening $(b_i,r_i)$, and the
aggregate OR proof has an extracted
opening $(t,r_T)$ with
$t\in\mathcal S_{m,b}$.
If $\sum_i b_i\neq b$ in $\F_2$,
then
\[
 t\not\equiv S:=\sum_i b_i\pmod P.
\]
\end{lemma}
\begin{proof}
The incorrect parity means
$S\not\equiv b\pmod2$,
whereas every $t\in\mathcal S_{m,b}$
satisfies $t\equiv b\pmod2$.
Thus $S\neq t$ as integers.
Both lie in $\{0,\ldots,m\}$ and
$P>m$, so they remain distinct
modulo $P$.
\end{proof}

\begin{theorem}[False parity reveals the setup logarithm]\label{thm:realize:parity-binding}
Under the hypotheses of
Lemma~\ref{lem:realize:parity-equivocation},
the bit openings and aggregate opening
explicitly determine the setup logarithm:
\[
 \boxed{
 \log_g h=(S-t)(r_T-R_T)^{-1}\pmod P,
 \quad R_T=\sum_i r_i.}
\]
Therefore a polynomial-time adversary
that yields extracted accepting proofs for
a false parity relation with non-negligible
probability contradicts the discrete-logarithm
assumption.
\end{theorem}
\begin{proof}
From the extracted bit openings,
Lemma~\ref{lem:realize:aggregate} gives
$C_T=g^S h^{R_T}$.
The extracted aggregate opening gives
$C_T=g^t h^{r_T}$.
Hence $g^{S-t}=h^{r_T-R_T}$.
The previous lemma gives
$S-t\ne0$ modulo $P$, so
$r_T-R_T$ cannot vanish modulo $P$:
otherwise the nonidentity element
$g^{S-t}$ would equal $1$.
The formula follows by exponent comparison
and inversion, exactly as in
Lemma~\ref{lem:realize:binding-reduction}.
Computing it from the alleged proof
extractor contradicts discrete-logarithm
hardness whenever such adversaries succeed
non-negligibly and extraction itself succeeds.
\end{proof}

\begin{theorem}[Perfect honest-verifier privacy of parity proofs]\label{thm:realize:parity-hvzk}
For a fixed public commitment vector and
public parity row, the full set of
honest-verifier bitness and aggregate
OR-proof transcripts has a
perfect simulator using only the
commitments and row coefficients.
The simulator requires no committed bits
or blinding scalars.
\end{theorem}
\begin{proof}
For each bitness proof and the aggregate
proof, use the witness-free perfect simulator
from Theorem~\ref{thm:realize:or-hvzk}.
The publicly determined aggregate
commitment is the product of the
given commitments, so the simulator can
construct all branch statements.
Honest provers choose independent
fresh OR-protocol randomness and
honest sigma verifiers choose independent
uniform challenges. Conditional on the
fixed public statements, each OR transcript
therefore has exactly its independently
simulated law. Multiplying these conditional
laws yields equality of the complete
joint transcript distribution.
\end{proof}

\begin{corollary}[Exact empty-row handling]\label{cor:realize:empty-row}
If a row contains no committed bits,
its parity constraint is true exactly when
its public target is $0$, and it can be checked
without any cryptographic proof.
\end{corollary}
\begin{proof}
The empty sum in $\F_2$ equals $0$.
Thus its equality to the target holds
precisely for target $0$.
The candidate sets are
$\mathcal S_{0,0}=\{0\}$ and
$\mathcal S_{0,1}=\varnothing$,
in agreement with direct checking.
\end{proof}

\begin{remark}[Extractor scope]
The parity binding theorem is an
\emph{extracted-witness} argument:
it assumes the standard knowledge
extractor succeeds on each accepted
Schnorr OR proof. Its explicit
computational error must be added
when this extractor is composed with
the outer ring-switching challenges.
It is not a claim of unconditional
information-theoretic soundness.
\end{remark}
